% Part: second-order-logic
% Chapter: syntax-and-semantics
% Section: introduction

\documentclass[../../../include/open-logic-section]{subfiles}

% 确保在编译主文件时，LaTeX 使用此工作目录。
% 本文件是二阶逻辑章节的部分内容。

\begin{document}

\olfileid{sol}{syn}{int}

\olsection{引言}

在一阶逻辑中，我们将给定语言的非逻辑符号（即常项符号、函数符号和谓词符号）与逻辑符号结合起来，以表达关于一阶结构的内容。这是通过满足的概念来实现的，它将结构~$\Struct{M}$、变元指派~$s$ 与公式~$!A$ 联系起来：$\Sat{M}{!A}[s]$ 成立，当且仅当 $!A$ 在将其常项符号、函数符号和谓词符号按 $\Struct{M}$ 解释，并将其自由变元按~$s$ 解释后，所表达的内容为真。等词~$\eq$ 的解释已内置于 $\Sat{M}{!A}[s]$ 的定义中，对~$\lforall$ 和 $\lexists$ 的解释亦然。前者总是被解释为结构论域~$\Domain{M}$ 上的恒等关系，而量词总是被解释为遍历整个论域。但关键在于，量化只允许针对论域的元素，因此只有对象变元能跟在量词之后。

在二阶逻辑中，语言和满足的定义均被扩展，以包含自由与约束的函数变元和谓词变元，以及对它们的量化。这些变元与函数符号和谓词符号的关系，就如同对象变元与常项符号的关系。它们在二阶逻辑的项与公式的构成中扮演相同的角色，对其量化也以类似方式处理。在\emph{标准}语义下，二阶量词遍历所有适当类型的对象（对函数变元，是 $\Domain{M}$ 到~$\Domain{M}$ 的 $n$ 元函数；对谓词变元，则是 $n$ 元关系）。例如，虽然 $\lforall[\Obj{v_0}][(\Obj{P^1_0}(\Obj{v_0}) \lor \lnot
  \Obj{P^1_0}(\Obj{v_0}))]$ 在一阶和二阶逻辑中都是一个公式，但在后者中我们还可以考虑 $\lforall[\Obj{V^1_0}][\lforall[\Obj{v_0}][(\Obj{V^1_0}(\Obj{v_0}) \lor
    \lnot \Obj{V^1_0}(\Obj{v_0}))]]$ 和 $\lexists[\Obj{V^1_0}][\lforall[\Obj{v_0}][(\Obj{V^1_0}(\Obj{v_0}) \lor
    \lnot \Obj{V^1_0}(\Obj{v_0}))]]$。由于它们不含自由变元，因而是二阶逻辑的语句。这里，$\Obj{V^1_0}$ 是一个二阶的一元谓词变元。$\Obj{V^1_0}$ 的容许解释与可赋予一元谓词符号（如 $\Obj{P^1_0}$）的解释相同，即 $\Domain{M}$ 的子集。因此，对它们的量化就等于说，对所有将 $\Domain{M}$ 的子集指派为 $\Obj{V^1_0}$ 的值的情况，$\lforall[\Obj{v_0}][(\Obj{V^1_0}(\Obj v_0) \lor \lnot
  \Obj{V^1_0}(v_0))]$ 都成立，或至少对一种情况成立。由于任何集合要么包含要么不包含一个给定的对象，这两个语句在任何结构中都为真。
\end{document}